\chapter*{MỞ ĐẦU}
\addcontentsline{toc}{chapter}{MỞ ĐẦU}

\section*{1. Tính cấp thiết của đề tài}
\addcontentsline{toc}{section}{1. Tính cấp thiết của đề tài}
Sự phát triển của các hệ thống thông tin vô tuyến thế hệ mới đặt ra yêu cầu đồng thời về tốc độ dữ liệu, độ tin cậy, độ trễ và khả năng phục vụ số lượng lớn thiết bị. Trong môi trường phổ tần hữu hạn, việc chỉ tăng công suất phát hoặc mở rộng băng thông không phải lúc nào cũng khả thi. Một hướng tiếp cận quan trọng là cải thiện cách mạng quản lý nhiễu và điều khiển môi trường truyền sóng. RSMA được đề xuất như một khung đa truy cập linh hoạt, trong đó thông điệp của mỗi người dùng được phân tách thành phần chung và phần riêng. Người nhận trước tiên giải mã luồng chung, thực hiện khử nhiễu liên tiếp, sau đó giải mã luồng riêng của mình. Cơ chế này cho phép điều chỉnh liên tục giữa hai thái cực: coi nhiễu là tạp âm và giải mã một phần nhiễu, từ đó tổng quát hóa nhiều chiến lược đa truy cập đã có \cite{Mao2022RSMAFundamentals,Mao2018RSMABridging,Clerckx2021NOMACritical}.

Song song với tiến bộ ở lớp đa truy cập, RIS và STAR--RIS cho phép thay đổi có chủ đích biên độ hoặc pha của sóng điện từ trên đường truyền gián tiếp. RIS phản xạ truyền thống chủ yếu phục vụ người dùng ở cùng một nửa không gian, trong khi STAR--RIS có khả năng đồng thời truyền và phản xạ, mở rộng vùng phủ sang hai phía của bề mặt \cite{Xu2021STARRIS,Mu2021STARAided,Khalid2022STARDesign}. Khi kết hợp RSMA với STAR--RIS, hệ thống có thêm bậc tự do để điều khiển nhiễu: RSMA xử lý nhiễu ở miền tín hiệu, còn STAR--RIS định hình kênh hiệu dụng. Tuy nhiên, lợi ích này đi kèm bài toán tối ưu có nhiều biến liên tục, gồm công suất luồng chung, công suất các luồng riêng, tỷ lệ phân bổ tốc độ chung, hệ số phân chia năng lượng và các pha truyền/phản xạ.

Các phương pháp tối ưu truyền thống như AO, SCA hoặc tìm kiếm theo codebook có ưu điểm là bám sát cấu trúc toán học của bài toán. Dù vậy, chúng thường phải lặp nhiều lần cho mỗi hiện thực kênh mới. Trong bối cảnh yêu cầu ra quyết định nhanh, chi phí tính toán trực tuyến có thể trở thành nút thắt. Học tăng cường sâu cung cấp một cách nhìn khác: thực hiện phần lớn chi phí trong giai đoạn huấn luyện, sau đó sử dụng mạng nơ-ron như một ánh xạ gần tức thời từ trạng thái kênh sang quyết định phân bổ tài nguyên. TD3 đặc biệt phù hợp với không gian hành động liên tục và được thiết kế để giảm sai lệch ước lượng giá trị nhờ hai critic, cập nhật chính sách trễ và làm trơn chính sách đích \cite{Fujimoto2018TD3}.

Đề tài vì vậy tập trung xây dựng một quy trình nghiên cứu thống nhất cho hệ thống SISO STAR--RIS hỗ trợ RSMA, trong đó TD3 là phương pháp học chính. Luận văn không chỉ quan tâm đến tổng tốc độ, mà còn xem xét tỷ lệ thỏa QoS, xác suất tất cả người dùng cùng thỏa QoS, mức vi phạm ràng buộc, độ trễ ra quyết định, khả năng mở rộng theo số phần tử STAR--RIS và độ tin cậy thống kê. Đây là cách tiếp cận cần thiết để tránh kết luận chỉ dựa trên một chỉ số hoặc một lần chạy.

\section*{2. Vấn đề và khoảng trống nghiên cứu}
\addcontentsline{toc}{section}{2. Vấn đề và khoảng trống nghiên cứu}
Các công trình về RSMA đã hình thành nền tảng lý thuyết tương đối đầy đủ cho việc phân chia thông điệp, thiết kế precoder và tối ưu tổng tốc độ dưới CSI không hoàn hảo \cite{Joudeh2016RSMASumRate,Li2020RSPractical}. Các nghiên cứu RIS đã phát triển mô hình tối ưu kết hợp giữa biến chủ động và thụ động, đồng thời chỉ ra các vấn đề thực tế về suy hao đường truyền, ước lượng kênh và giới hạn phần cứng \cite{Wu2019IRSBeamforming,Huang2019RISEnergyEfficiency,DiRenzo2020SmartRadio,Bjornson2020RISMyths,Wu2021IRSTutorial}. Với STAR--RIS, ba giao thức energy splitting, mode switching và time switching được nghiên cứu rộng rãi; mô hình pha ghép cũng cho thấy giả định pha truyền và pha phản xạ độc lập có thể là lý tưởng hóa đối với phần cứng thụ động \cite{Mu2021STARAided,Liu2021STARCoupledPhase,Wang2022STAROptimization}.

Một số công trình gần đây đã kết hợp trực tiếp STAR--RIS và RSMA. Meng và cộng sự sử dụng PPO để tối đa hóa tổng tốc độ trong mạng STAR--RIS--RSMA \cite{Meng2024PPOSTARRSMA}; Ge và cộng sự nghiên cứu chiến lược rate-splitting cho massive MIMO với tối ưu chung \cite{Ge2024RSMAMassiveMIMO}; Chang và cộng sự sử dụng AO và penalty-SCA trong bài toán truyền thông bí mật \cite{Chang2024CovertSTARRSMA}. Những nghiên cứu này chứng minh tiềm năng của việc kết hợp hai công nghệ, nhưng mô hình anten, mục tiêu, ràng buộc và giao thức đánh giá khác nhau đáng kể. Do đó, không thể trực tiếp chuyển công thức hoặc kết quả của một công trình sang mô hình SISO của luận văn.

Khoảng trống mà luận văn hướng đến gồm bốn điểm. Thứ nhất, cần một mô hình SISO tối giản nhưng nhất quán, trong đó tất cả phương pháp sử dụng cùng phương trình kênh và cùng bộ tính tốc độ RSMA. Thứ hai, cần một không gian hành động khả thi về mặt vật lý, tránh tình trạng actor sinh ra nghiệm vi phạm tổng công suất hoặc vi phạm quy tắc phân chia năng lượng. Thứ ba, cần quy trình đánh giá khóa train/validation/test để loại bỏ rò rỉ dữ liệu và lựa chọn checkpoint trên tập test. Thứ tư, cần diễn giải hiệu năng dưới góc nhìn đánh đổi chất lượng--độ trễ thay vì chỉ khẳng định một thuật toán ``tốt nhất'' theo tổng tốc độ.

\section*{3. Mục tiêu nghiên cứu}
\addcontentsline{toc}{section}{3. Mục tiêu nghiên cứu}
\subsection*{3.1. Mục tiêu tổng quát}
Mục tiêu tổng quát của luận văn là xây dựng và đánh giá phương pháp TD3 đơn tác tử để tối ưu đồng thời tài nguyên RSMA và cấu hình STAR--RIS trong hệ thống truyền xuống SISO, qua đó đạt sự cân bằng giữa tổng tốc độ truyền, mức độ đáp ứng QoS và độ trễ ra quyết định.

\subsection*{3.2. Mục tiêu cụ thể}
Các mục tiêu cụ thể gồm:
\begin{enumerate}[label=\arabic*)]
  \item Xây dựng mô hình toán học cho hệ thống SISO STAR--RIS energy-splitting hỗ trợ RSMA với người dùng ở cả phía truyền và phía phản xạ.
  \item Xây dựng bài toán tối ưu đồng thời công suất luồng chung/riêng, tỷ lệ phân bổ tốc độ chung, hệ số phân chia năng lượng và pha truyền/phản xạ dưới ràng buộc tổng công suất và QoS.
  \item Mô hình hóa bài toán thành một quá trình quyết định theo ngữ cảnh và phát triển thuật toán TD3 cho không gian hành động liên tục có số chiều tăng theo số phần tử STAR--RIS.
  \item Xây dựng các baseline truyền thống và học tăng cường trong cùng một môi trường vật lý; thực hiện ablation để đánh giá vai trò của STAR--RIS và chính sách phân bổ công suất.
  \item Đánh giá bằng ScenarioBank khóa, nhiều seed, thống kê ghép cặp, khoảng tin cậy, kiểm định giả thuyết và benchmark độ trễ CPU một luồng.
\end{enumerate}

\section*{4. Câu hỏi nghiên cứu}
\addcontentsline{toc}{section}{4. Câu hỏi nghiên cứu}
\begin{researchquestion}
\textbf{RQ1:} TD3 có học được ánh xạ từ trạng thái kênh sang phương án phân bổ tài nguyên khả thi và duy trì QoS ổn định trên các kịch bản chưa quan sát hay không?
\end{researchquestion}
\begin{researchquestion}
\textbf{RQ2:} Chất lượng nghiệm của TD3, xét theo tổng tốc độ và vi phạm QoS, khác như thế nào so với AO--SCA, AO--Grid và AnalyticalRIS?
\end{researchquestion}
\begin{researchquestion}
\textbf{RQ3:} Khi số phần tử STAR--RIS tăng từ 16 đến 128, tổng tốc độ, độ tin cậy QoS và độ trễ của các phương pháp thay đổi theo xu hướng nào?
\end{researchquestion}
\begin{researchquestion}
\textbf{RQ4:} TD3 tạo ra điểm đánh đổi chất lượng--độ trễ ra sao so với các solver lặp truyền thống?
\end{researchquestion}

\section*{5. Đối tượng và phạm vi nghiên cứu}
\addcontentsline{toc}{section}{5. Đối tượng và phạm vi nghiên cứu}
Đối tượng nghiên cứu gồm: hệ thống truyền xuống SISO STAR--RIS hỗ trợ RSMA; mô hình energy-splitting của STAR--RIS; bài toán phân bổ công suất và tốc độ chung; thuật toán TD3 cùng các baseline; và quy trình đánh giá thực nghiệm có tính tái lập.

Phạm vi mô hình được giới hạn ở một trạm gốc SISO, bốn người dùng SISO, một STAR--RIS thụ động và các kích thước $N\in\{16,32,64,96,128\}$. Người dùng được gán cố định vào phía truyền hoặc phía phản xạ. CSI được cung cấp cho bộ ra quyết định thông qua observation. Kênh được sinh tổng hợp theo phân bố Gaussian phức với các hệ số tỷ lệ cố định; chưa xét vị trí hình học, suy hao đường truyền theo khoảng cách, kênh Rician, di động, tương quan thời gian, đa ô, active STAR--RIS, lỗi phần cứng, sai số ước lượng CSI hoặc thử nghiệm phần cứng.

Phạm vi thuật toán tập trung vào TD3. DDPG và PPO được trình bày như baseline học tăng cường trong framework; tuy nhiên, các tuyên bố định lượng cuối cùng của luận văn chỉ sử dụng những kết quả đã có đầy đủ kiểm toán và provenance trong bundle thực nghiệm. AO--SCA được xem là baseline tối ưu cục bộ, không phải nghiệm tối ưu toàn cục hay upper bound.

\section*{6. Phương pháp nghiên cứu}
\addcontentsline{toc}{section}{6. Phương pháp nghiên cứu}
Luận văn sử dụng kết hợp các phương pháp sau:
\begin{itemize}
  \item \textbf{Tổng hợp tài liệu:} phân tích các công trình nền tảng về RSMA, RIS, STAR--RIS, TD3 và SCA từ thư viện PDF đã kiểm tra trong repository.
  \item \textbf{Mô hình hóa toán học:} xây dựng kênh hiệu dụng, SINR, tốc độ chung/riêng, hàm mục tiêu và tập ràng buộc.
  \item \textbf{Thiết kế thuật toán:} ánh xạ observation và action sang không gian vật lý; thiết kế reward phạt QoS; áp dụng dual ascent chiếu để điều chỉnh hệ số phạt.
  \item \textbf{Mô phỏng định lượng:} triển khai bằng Python và PyTorch, huấn luyện nhiều seed trên các ScenarioBank độc lập.
  \item \textbf{Đánh giá thống kê:} tính trung bình, độ lệch chuẩn, khoảng tin cậy 95\%, kiểm định $t$ ghép cặp, Wilcoxon ghép cặp, hiệu chỉnh Holm và effect size.
  \item \textbf{Đánh giá hệ thống:} đo inference, tính metric và end-to-end latency trong chế độ CPU một luồng.
\end{itemize}

\section*{7. Đóng góp của luận văn}
\addcontentsline{toc}{section}{7. Đóng góp của luận văn}
Các đóng góp chính gồm:
\begin{enumerate}[label=\arabic*)]
  \item Một mô hình SISO STAR--RIS--RSMA thống nhất, trong đó các hệ số STAR--RIS, công suất và tỷ lệ tốc độ được giải mã từ cùng một vector hành động khả thi.
  \item Một bộ tối ưu TD3 có cơ chế hai critic, target-policy smoothing, delayed update, normalization theo khối, action projection và điều khiển dual cho QoS.
  \item Một quy trình ScenarioBank khóa train/validation/test, checkpoint selection theo ưu tiên tính khả thi, đánh giá test không thăm dò và lưu provenance bằng checksum/config hash/Git commit.
  \item Một đánh giá đa chiều về sum-rate, QoS, violation, scalability, ablation, latency và ý nghĩa thống kê.
  \item Một diễn giải khoa học thận trọng: TD3 được định vị là giải pháp cân bằng chất lượng--độ trễ có QoS cao, trong khi AO--SCA đạt tổng tốc độ cao hơn trong thí nghiệm đã thực hiện.
\end{enumerate}

\section*{8. Kết cấu luận văn}
\addcontentsline{toc}{section}{8. Kết cấu luận văn}
Ngoài phần Mở đầu, Kết luận, Tài liệu tham khảo và Phụ lục, luận văn gồm năm chương. Chương 1 tổng quan các hướng nghiên cứu và xác định khoảng trống. Chương 2 trình bày cơ sở lý thuyết về RSMA, STAR--RIS, học tăng cường và TD3. Chương 3 xây dựng mô hình hệ thống và bài toán tối ưu. Chương 4 mô tả chi tiết phương pháp TD3, các baseline và quy trình thực nghiệm. Chương 5 phân tích kết quả mô phỏng, độ trễ, kiểm định thống kê, ablation và các mối đe dọa đến tính hợp lệ của kết luận.
